\documentclass[10pt,a4paper]{article}
\usepackage[margin=24mm]{geometry}
\usepackage{amsmath,amssymb,booktabs}
\usepackage[T1]{fontenc}
\usepackage{lmodern}
\usepackage[hidelinks]{hyperref}
\setlength{\parindent}{0pt}
\setlength{\parskip}{4pt}
\newcommand{\F}{\mathbb F}
\newcommand{\wt}{\operatorname{wt}}
\title{Problem 00000001082: the full-dimension MRD boundary}
\author{\texttt{lizaixi01}}
\date{24 October 2026}
\begin{document}
\begin{center}
{\Large\bfseries Problem 00000001082:\\the full-dimension MRD boundary}\par
\smallskip
\texttt{lizaixi01}\quad -- \quad 24 October 2026
\end{center}

\section*{Statement and scope}
The printed conjecture gives, for fixed $(m,n,k)$, the number of MRD-code
equivalence classes as
\[
  q^{(k-1)(n-k+1)}\bigl(1+O(q^{-1})\bigr).
\]
Use the standard convention that a linear rank-metric code is an
$\F_{q^m}$-linear subspace of $\F_{q^m}^{\,n}$ and $k$ is its dimension
over that extension field. Its dimension over $\F_q$ is $mk$, not $k$.
We take
\[
  (m,n,k)=(3,2,2).
\]
The source prints no restriction $k<n$ or minimum distance greater than one.
This choice also satisfies the strict $n<m$ in Morrison's Definition II.11
\cite{Morrison}. The formal result concerns the asserted count of MRD codes
in this standard linear model. It makes no separate Gabidulin-only
classification claim and does not replace $k$ by a matrix-space dimension.

It is enough to let $q=p$ range over all primes: an asymptotic formula for
all prime powers must hold on this unbounded subfamily. We will derive the
exact count $1$ for every such $p$, and prove that it cannot have the
displayed main term $p$ and relative error $O(p^{-1})$.

\section*{The actual code and its distance}
Fix any prime $p$. Let $K=\F_p$, let $L=\F_{p^3}$, and put $V=L^2$.
The Lean construction uses the native field
\texttt{GaloisField p 3}, a splitting field over \texttt{ZMod p}; its
actual extension degree is $3$ and its cardinality is $p^3$.
For $v=(v_1,v_2)\in V$ define the standard rank weight and distance by
\[
 \wt(v)=\dim_K\langle v_1,v_2\rangle_K,
 \qquad d_R(x,y)=\wt(x-y).
\]
This is Morrison's Definition II.8 and Sheekey's vector model
\cite{Morrison,Sheekey}. For a code $C$ the minimum is taken over
\emph{all} pairs $x,y\in C$ with $x\ne y$.

Every nonzero vector has positive rank weight: if its coordinate span had
dimension zero, that span would be zero and both coordinates would vanish.
The actual vector $u=(1,0)$ has coordinate span $K\cdot1$ of dimension
one. Hence the full code $C=V$ has minimum distance exactly $1$, attained
at $(u,0)$. Its distance set is nonempty; no empty-infimum convention is used.
Its actual size is
\[
 |C|=|L|^{\dim_L C}=(p^3)^2=p^6
     =p^{\,3(2-d_R(C)+1)}.
\]
This is precisely equality in the rank-metric Singleton bound, so $C$ is
MRD \cite{Sheekey}. Both the minimum and the cardinality are calculated
from the actual code, rather than inserted as its defining values.

\section*{The complete equivalence-class count}
Let $D$ be \emph{any} $L$-linear code of length $2$ and dimension $2$.
Since $D\subseteq V$ and $\dim_L D=\dim_L V=2$, finite-dimensional linear
algebra gives $D=V$. Thus the complete universe of $[2,2]$ linear MRD
codes consists of the single actual subspace $V$.

We use the standard linear rank-equivalence from Morrison's
Definition III.4: $D$ and $E$ are equivalent if an invertible $L$-linear
map $f:V\to V$ preserves the actual rank weight of every vector and
maps $D$ onto $E$. Identity, inverse and composition verify that this is
an equivalence relation. The Lean quotient uses this relation, including
the actual image-of-code condition. It is nonempty and a subsingleton,
so its finite cardinality is exactly
\[
  A(p)=1.
\]
Uniqueness was proved before forming the quotient. Consequently allowing
additional standard semilinear equivalences cannot change this one-class
count; no arbitrary relation or chosen list of representative examples is
being used to obtain the result.

\section*{Failure of the claimed asymptotic}
For the fixed parameters, the printed exponent is
$(2-1)(2-2+1)=1$. The true relative error is therefore
\[
 R(p)=\frac{A(p)}p-1=\frac1p-1,
 \qquad |R(p)|=\frac{p-1}p.
\]
For every real constant $B$ and every integer cutoff $N$, Euclid's theorem
provides a prime $p>\max\{N,B+1\}$. At that actual prime,
\[
  |R(p)|=\frac{p-1}p>\frac Bp.
\]
Thus $R(p)$ is not $O(p^{-1})$, with all constants and eventual cutoffs
quantified. More explicitly, if $A(p)=p(1+\epsilon(p))$ eventually and
$\epsilon(p)=O(p^{-1})$, division by the nonzero $p$ forces
$\epsilon(p)=R(p)$ eventually, contradicting the preceding inequality.
The formal proof uses native \texttt{Asymptotics.IsBigO} on the ordered
type of all primes with its \texttt{atTop} filter, as well as the explicit
arbitrarily-large-prime witnesses.

This disproves the unrestricted asymptotic clause as printed. It makes
no assertion about variants excluding $k=n$, nontrivial minimum-distance
families, or the unspecified cell-decomposition clause.

\section*{Formalization and reproducibility}
The complete proof is in \texttt{lean/Main.lean}, namespace
\texttt{TLMC1082}. Its principal declarations are:
\begin{center}
\begin{tabular}{ll}
\toprule
Bridge & Declaration\\
\midrule
Actual finite-field parameters & \texttt{actual\_parameters}\\
True minimum and attainment & \texttt{full\_minimumDistance}\\
Actual Singleton equality & \texttt{full\_isMRD}\\
Entire code classification & \texttt{code\_eq\_top}\\
Standard equivalence quotient & \texttt{classCount\_one}\\
All constants and cutoffs fail & \texttt{arbitrarily\_large\_prime\_failure}\\
Native asymptotic impossibility & \texttt{printed\_asymptotic\_false}\\
Main result & \texttt{complete\_counterexample}\\
\bottomrule
\end{tabular}
\end{center}
The project pins Lean~4.33.0 and Mathlib commit\par
\texttt{db584cd6d46c92f209a44c0f1c829460d327499d}.\par
First run \texttt{lean --version}; then \texttt{lean/verify.sh} runs the
root build, warning-as-error namespace axioms audit and official full
fresh-kernel replay. Every owned declaration, including generated names,
is enumerated and checked with ordinary \texttt{\#print axioms}.
The allowed foundational axioms are only \texttt{propext},
\texttt{Classical.choice}, and \texttt{Quot.sound}.

\begin{thebibliography}{9}
\bibitem{Morrison}
K. Morrison, \emph{Equivalence for Rank-metric and Matrix Codes and
Automorphism Groups of Gabidulin Codes}, arXiv:1304.0501v1,
Definitions II.8, II.11 and III.4, printed pp.~4, 5 and 8.
\url{https://arxiv.org/abs/1304.0501v1}.
\bibitem{Sheekey}
J. Sheekey, \emph{MRD Codes: Constructions and Connections},
arXiv:1904.05813v1 (2019), Sections 1.1--1.2 and Definition 3,
PDF pp.~2--3. \url{https://arxiv.org/abs/1904.05813v1}.
\end{thebibliography}
\end{document}
